% Generated draft: requires mathematical and source review.
\documentclass[11pt]{article}
\usepackage[margin=25mm]{geometry}
\usepackage{fontspec}
\setmainfont{Latin Modern Roman}
\usepackage{amsmath,amssymb,mathtools,unicode-math}
\setmathfont{Latin Modern Math}
\usepackage{enumitem,needspace,xurl,hyperref}
\setlength{\parindent}{0pt}
\setlength{\parskip}{5pt}
\setlength{\emergencystretch}{3em}
\newcommand{\N}{\mathbb{N}}
\newcommand{\Z}{\mathbb{Z}}
\newcommand{\Q}{\mathbb{Q}}
\newcommand{\R}{\mathbb{R}}
\newcommand{\C}{\mathbb{C}}
\newcommand{\F}{\mathbb{F}}
\newcommand{\A}{\mathbb{A}}
\newcommand{\PP}{\mathbb{P}}
\newcommand{\E}{\mathbb{E}}
\newcommand{\eps}{\varepsilon}
\newcommand{\dd}{\,\mathrm{d}}
\begin{document}
\section*{UnsolvedMath ID: 6200074}
\textbf{Source problem code:} catalogue entry
\subsection{Definitions and mathematical statement}
A finitely generated group $G$, with a word metric $d$, is lacunary hyperbolic in the standard sense if for some sequence $s_i\to\infty$ and some nonprincipal ultrafilter $\omega$, the asymptotic cone
\[
\operatorname{Cone}_\omega(G,(s_i))=\lim_\omega(G,s_i^{-1}d,e)
\]
is an $\mathbb R$-tree. An $\mathbb R$-tree is a geodesic metric space in which the unique arc between any two points is a geodesic segment. The condition requires an $\mathbb R$-tree cone, not that every asymptotic cone be an $\mathbb R$-tree.

The source asks three questions for this class:
\begin{enumerate}
\item Develop a meaningful structure theory for lacunary hyperbolic groups.
\item Define a useful boundary for such groups, with appropriate properties still to be specified.
\item Is it true that for every such $G$, either $\operatorname{Out}(G)$ is finite or $G$ has a nontrivial splitting over a virtually cyclic subgroup?
\end{enumerate}
Here $\operatorname{Out}(G)=\operatorname{Aut}(G)/\operatorname{Inn}(G)$. A virtually cyclic group has a cyclic finite-index subgroup. A nontrivial splitting can be expressed as an action on a Bass--Serre tree without a global fixed point, with edge stabilizers virtually cyclic.

The source ends with an apparent footnote saying such groups are never finitely presented. Under the standard definition above, ordinary hyperbolic groups are included, so that blanket restriction cannot be adopted as part of the definition. The intended footnote may concern only nonhyperbolic lacunary hyperbolic groups, but its context is not supplied. The class intended in the three questions therefore needs checking.
\subsection{Short English statement}
Develop structure and boundary theories for lacunary hyperbolic groups. Decide whether infinite outer automorphism group forces a splitting over a virtually cyclic subgroup, and clarify whether only nonhyperbolic examples are intended.
\subsection{Sources}
\begin{itemize}
\item Misha Kapovich, \emph{Problems on Boundaries of Groups and Kleinian Groups} (2005), Problem 74, PDF page 19, \url{https://www.math.ucdavis.edu/~kapovich/EPR/problems.pdf}; accessed 2026-07-29, PDF-text-normalized extraction.
\item The supplied summary attributes introduction of lacunary hyperbolic groups to Ol'shanskii--Osin--Sapir and mentions boundary work described as ``boundaries of lacunary hyperbolic groups'' by Kar/Weidmann; no publication identifiers are supplied.
\item Local registry \texttt{data/unsolvedmath/problems.json}, permanent UnsolvedMath ID $6200074$, source code catalogue entry, \emph{Boundaries of Groups and Kleinian Groups --- Problem 74}. Dataset attribution: Ulam AI and the UnsolvedMath Contributors.
\end{itemize}
\end{document}
